\chapter{Templates \& Classifier Specifications}
\label{ch:templates}


Reusable artifacts for putting the framework into practice.

\section{EPG Constraint Library Template}
\label{sec:templates-C-1}


\subsection{Library Structure}
\label{sec:templates-C-1-1}


A constraint library package contains:

\textbf{Library metadata:}
\begin{itemize}
  \item Name and version (semantic versioning: major.minor.patch)
  \item Target industry and regulatory basis (e.g., ``HIPAA Privacy Rule \S{}164.502-\S{}164.514")
  \item Maintainer organization and contact
  \item Last review date and next scheduled review
  \item Applicable governance levels (recommended)
\end{itemize}


\textbf{Constraint entries:} each containing:
\begin{itemize}
  \item Deontic modality (O or F)
  \item Atomic property identifier from the controlled vocabulary
  \item Recommended governance level (enterprise, department, or project)
  \item Domain assignment
  \item Evaluability class (C\_\allowbreak{}m or C\_\allowbreak{}p)
  \item Rationale and regulatory reference
\end{itemize}


\textbf{Pre-attested decompositions:} for each semantic intent underlying the constraints:
\begin{itemize}
  \item The intent record (C\_\allowbreak{}s) with full semantic description
  \item The mechanical components (C\_\allowbreak{}m) with implementation specifications
  \item The procedural components (C\_\allowbreak{}p) with attestation requirements
  \item The decomposition attestation record (who attested, when, qualifications)
\end{itemize}


\textbf{Known residual gaps:} documented per constraint where soundness cannot be fully demonstrated, including compensating control recommendations.

\textbf{Regression test suite:} known-bad inputs that mechanical checks must catch, plus expected evaluation outcomes. This suite is the minimum test set that must pass during local verification (\S{}9.4).

\subsection{Library Lifecycle}
\label{sec:templates-C-1-2}


Libraries follow semantic versioning:

\begin{itemize}
  \item \textbf{Major version} (e.g., 1.0 $\to$ 2.0): new constraints added, constraints removed, or constraint semantics materially changed. Requires full re-attestation by adopting organizations.
  \item \textbf{Minor version} (e.g., 1.0 $\to$ 1.1): modified decompositions, updated mechanical checks, expanded test suites. Requires regression test re-execution by adopting organizations.
  \item \textbf{Patch version} (e.g., 1.0.0 $\to$ 1.0.1): documentation updates, test case additions, editorial corrections. No action required by adopting organizations beyond acknowledgment.
\end{itemize}


Each version requires domain expert attestation from the library maintainer. Organizations that have adopted a library version receive update notifications when new versions are released, including a changelog and impact assessment.

\subsection{Example: HIPAA Prompt Constraints v1.0 (Skeleton)}
\label{sec:templates-C-1-3}


\textbf{Library metadata:}
\begin{itemize}
  \item Name: HIPAA Prompt Constraints
  \item Version: 1.0.0
  \item Industry: Healthcare / Life Sciences
  \item Regulatory basis: HIPAA Privacy Rule (45 CFR \S{}164.502-\S{}164.514), HIPAA Security Rule (45 CFR \S{}164.302-\S{}164.318)
  \item Recommended level: Enterprise
\end{itemize}


\textbf{Constraint entries (8 constraints):}

\begin{table}[ht]
  \centering\footnotesize
  \begin{tabular}{@{}>{\raggedright\arraybackslash}p{1.12cm}>{\raggedright\arraybackslash}p{1.12cm}>{\raggedright\arraybackslash}p{4.11cm}>{\raggedright\arraybackslash}p{1.19cm}>{\raggedright\arraybackslash}p{6.11cm}@{}}
  \toprule
  \# & Modality & Property & Class & Regulatory Reference \\
  \midrule
  1 & F & phi\_\allowbreak{}disclosure\_\allowbreak{}in\_\allowbreak{}output & C\_\allowbreak{}m & \S{}164.502(a) --- Minimum Necessary \\
  2 & O & hipaa\_\allowbreak{}disclaimer\_\allowbreak{}present & C\_\allowbreak{}m & \S{}164.520 --- Notice of Privacy Practices \\
  3 & F & pii\_\allowbreak{}in\_\allowbreak{}system\_\allowbreak{}logs & C\_\allowbreak{}m & \S{}164.312(a) --- Access Control \\
  4 & O & data\_\allowbreak{}minimization\_\allowbreak{}instruction & C\_\allowbreak{}m & \S{}164.502(b) --- Minimum Necessary Standard \\
  5 & F & medical\_\allowbreak{}diagnosis\_\allowbreak{}instruction & C\_\allowbreak{}m + C\_\allowbreak{}p & \S{}164.502 --- Uses and Disclosures \\
  6 & O & patient\_\allowbreak{}redirect\_\allowbreak{}to\_\allowbreak{}provider & C\_\allowbreak{}m & \S{}164.502 --- Treatment Exception \\
  7 & F & phi\_\allowbreak{}in\_\allowbreak{}test\_\allowbreak{}environments & C\_\allowbreak{}m & \S{}164.308(a)(4) --- Information Access Management \\
  8 & O & breach\_\allowbreak{}notification\_\allowbreak{}instruction & C\_\allowbreak{}p & \S{}164.404 --- Notification to Individuals \\
  \bottomrule
  \end{tabular}
\end{table}


\textbf{Decomposition example (Constraint \#1: F(phi\_\allowbreak{}disclosure\_\allowbreak{}in\_\allowbreak{}output)):}

Intent record: ``The AI system must not include protected health information in its outputs unless the disclosure is explicitly authorized under a HIPAA-permitted use or disclosure."

Mechanical components:
\begin{itemize}
  \item NER-based PHI detection in prompt instructions (names, dates, MRNs, SSNs, addresses, phone numbers)
  \item Pattern-based detection of prompt instructions that request PHI inclusion
  \item Required presence of explicit ``Do not include patient-identifying information" instruction
\end{itemize}


Procedural components:
\begin{itemize}
  \item Privacy officer review and attestation that the prompt's PHI protections are adequate for the intended use case
  \item Annual re-attestation aligned with HIPAA risk assessment cycle
\end{itemize}


Residual gap: mechanical NER detection covers approximately 90\% of standard PHI patterns. Novel identifiers, contextual PHI (e.g., ``the patient in room 3"), and inference-based re-identification are not fully captured. Compensating control: 10\% HITL review of outputs by trained privacy staff.

\textbf{Regression test suite (partial):}
\begin{itemize}
  \item 50 known-PHI test prompts containing explicit patient identifiers $\to$ mechanical checks must flag all 50
  \item 20 boundary-case prompts with contextual identifiers $\to$ mechanical checks should flag at least 15
  \item 30 clean prompts with no PHI $\to$ mechanical checks must pass all 30 (false positive baseline)
\end{itemize}






\section{ROC Output Classifier Specifications}
\label{sec:templates-C-2}


\subsection{Classifier Governance Profile Template}
\label{sec:templates-C-2-1}


Every classifier deployed in ROC must have a documented governance profile:

\small
\begin{longtable}{@{}>{\raggedright\arraybackslash}p{2.48cm}>{\raggedright\arraybackslash}p{12.22cm}@{}}
  \toprule
  Field & Description \\
  \midrule
\endfirsthead
  \toprule
  Field & Description \\
  \midrule
\endhead
  \bottomrule
\endlastfoot
  Classifier ID & Unique identifier \\
  Version & Semantic version (major.minor.patch) \\
  Purpose & What threat/constraint it evaluates \\
  Output constraint class & O\_\allowbreak{}c or component of O\_\allowbreak{}x \\
  Training data summary & Source, size, composition (no PII in the profile) \\
  Benchmark metrics & Precision, recall, F1 at configured threshold \\
  Known limitations & Classes of inputs where performance degrades \\
  Expected Calibration Error (ECE) & Calibration metric for the current version --- required for threshold governance. A threshold is meaningless without knowing the classifier's calibration curve. \\
  Reliability diagram & Visual calibration curve for the current version \\
  Threshold ($\tau$) & Current governance threshold, with risk-tier variants. Only governable when ECE is known. \\
  Threshold owner & Which governance domain owns the threshold \\
  Validation cadence & How often retested against updated benchmarks \\
  Last validation date & When last validated \\
  Red-team results & Summary of last adversarial testing \\
  Retraining trigger & Conditions that mandate classifier retraining \\
\end{longtable}
\normalsize


\subsection{Minimum Classifier Requirements by Threat}
\label{sec:templates-C-2-2}


\begin{table}[ht]
  \centering\small
  \begin{tabular}{@{}>{\raggedright\arraybackslash}p{2.66cm}>{\raggedright\arraybackslash}p{5.45cm}>{\raggedright\arraybackslash}p{6.24cm}@{}}
  \toprule
  Threat & Minimum Classifier Type & Recommended Benchmark \\
  \midrule
  T7: Prompt injection & Sequence classifier on output patterns & Injection benchmark dataset (updated quarterly) \\
  T8: Jailbreak & Safety boundary classifier & Red-team jailbreak corpus (updated quarterly) \\
  T10: PII/PHI & NER + regex composite & HIPAA PHI pattern corpus, PCI test data \\
  T10: Memorization & Similarity scorer against training corpus & Memorization benchmark (model-specific) \\
  T11: Cross-tenant & Tenant-specific terminology detector & Tenant vocabulary cross-reference \\
  \bottomrule
  \end{tabular}
\end{table}


\subsection{Threshold Calibration Guidance}
\label{sec:templates-C-2-3}


Threshold selection involves a precision-recall tradeoff with regulatory implications:

\textbf{Conservative (low $\tau$, \textasciitilde{}0.5-0.7):} More flags, fewer misses. Higher false-positive rate increases review burden. Appropriate for Critical-tier PHI/safety constraints where a missed violation has severe regulatory or safety consequences.

\textbf{Balanced (medium $\tau$, \textasciitilde{}0.7-0.85):} Moderate flags and misses. Appropriate for Standard-tier constraints where the cost of false positives and false negatives is roughly symmetric.

\textbf{Permissive (high $\tau$, \textasciitilde{}0.85-0.95):} Fewer flags, more misses. Lower review burden but higher miss rate. Appropriate for Low-tier or informational constraints where violations are not safety-critical.

Threshold calibration is not a one-time event. As models change, user populations shift, and adversarial techniques evolve, the operating point on the precision-recall curve shifts. Periodic recalibration (aligned with classifier validation cadence) ensures thresholds remain appropriate.




